\documentclass[10pt]{article}

\usepackage[margin=1.8cm]{geometry}
\usepackage{enumitem}
\usepackage{titlesec}

\setlist[itemize]{
    itemsep=6pt,
    topsep=4pt,
    parsep=0pt,
    partopsep=0pt,
    leftmargin=1.4em
}

\titlespacing*{\section}{0pt}{10pt}{6pt}
\titlespacing*{\subsection}{0pt}{8pt}{4pt}
\titlespacing*{\paragraph}{0pt}{4pt}{4pt}

\titleformat{\section}
  {\normalfont\large\bfseries}
  {}
  {0pt}
  {\titlerule\vspace{0.5ex}}

\setlength{\parindent}{0pt}
\setlength{\parskip}{4pt}

\newcommand{\mainsection}[1]{%
  \begin{center}
    \Large\bfseries #1
  \end{center}
  \vspace{6pt}
}
\newcommand{\clabel}[1]{{\large\bfseries #1}}

\title{IHTC Metamodel}
\author{Maksym Byelko}
\date{July 2026}

\begin{document}

\maketitle
\mainsection{METAMODEL SUMMARY}

\section*{Overview}
This metamodel captures the current IHTC work in MDEO for two connected optimisation layers: patient scheduling and nurse assignment. The scheduling side decides which patients are admitted, on what day, in which room, and in which operating theatre. The nurse-assignment side then represents room-shift demand and assigns nurses to those shifts. The model also includes day-based availability objects so that transformations can update and validate state over time.

\section*{Enumerations}
\begin{itemize}
    \item \clabel{Gender.}\\
    \textbf{Purpose:} stores patient sex as \texttt{M} or \texttt{F}.

    \item \clabel{AgeGroup.}\\
    \textbf{Purpose:} stores patient age category as \texttt{BABY}, \texttt{CHILD}, \texttt{YOUNG}, \texttt{ADULT}, \texttt{ELDERLY}, or \texttt{EMPTY}.\\
    \textbf{Nuance:} \texttt{EMPTY} is also used in \texttt{RoomAvailability} as a sentinel value for a room-day that currently has no occupants.
\end{itemize}

\section*{Root Container}
\begin{itemize}
    \item \clabel{HospitalInstance.}\\
    \textbf{Purpose:} root object that collects the full optimisation state for one hospital instance.
    \begin{itemize}
        \item \textbf{Attribute:} \texttt{decisionHorizon : int}
        \item \textbf{Link:} \texttt{patients[0..*] -> Patient}
        \item \textbf{Link:} \texttt{operatingtheatres[0..*] -> OperatingTheatre}
        \item \textbf{Link:} \texttt{surgeons[0..*] -> Surgeon}
        \item \textbf{Link:} \texttt{rooms[0..*] -> Room}
        \item \textbf{Link:} \texttt{nurses[0..*] -> Nurse}
        \item \textbf{Link:} \texttt{surgeonAvailabilities[0..*] -> SurgeonAvailability}
        \item \textbf{Link:} \texttt{operatingTheatreAvailabilities[0..*] -> OperatingTheatreAvailability}
        \item \textbf{Link:} \texttt{roomAvailabilities[0..*] -> RoomAvailability}
        \item \textbf{Link:} \texttt{hospitalisationShifts[0..*] -> HospitalisationShift}
        \item \textbf{Link:} \texttt{nurseWorkingShifts[0..*] -> NurseWorkingShift}
        \item \textbf{Link:} \texttt{roomShiftAssignments[0..*] -> RoomShiftAssignment}
        \item \textbf{Link:} \texttt{deletedAdmissionsTrackers[0..*] -> DeletedAdmissionsTracker}
    \end{itemize}
    \textbf{Nuance:} the \texttt{decisionHorizon} attribute defines the number of days in scope for the planning problem.
\end{itemize}

\section*{Patient Scheduling Core}
\begin{itemize}
    \item \clabel{Patient.}\\
    \textbf{Purpose:} represents a patient that may or may not be admitted.
    \begin{itemize}
        \item \textbf{Attribute:} \texttt{id : int}
        \item \textbf{Attribute:} \texttt{isMandatory : boolean}
        \item \textbf{Attribute:} \texttt{isScheduled : boolean}
        \item \textbf{Attribute:} \texttt{dueDate : int}
        \item \textbf{Attribute:} \texttt{releaseDate : int}
        \item \textbf{Attribute:} \texttt{ageGroup : AgeGroup}
        \item \textbf{Attribute:} \texttt{surgeryDuration : int}
        \item \textbf{Attribute:} \texttt{gender : Gender}
        \item \textbf{Attribute:} \texttt{stayLength : int}
        \item \textbf{Link:} \texttt{assignedSurgeonId[1..1] -> Surgeon}
        \item \textbf{Link:} \texttt{incompatibleRooms[0..*] -> Room}
        \item \textbf{Link:} \texttt{dayDemand[1..*] -> PatientDayDemand}
    \end{itemize}
    \textbf{Nuance:} \texttt{isScheduled} is currently used as a practical modelling flag to simplify matching and transformation logic.

    \item \clabel{Admission.}\\
    \textbf{Purpose:} records the actual scheduling decision for a patient.
    \begin{itemize}
        \item \textbf{Attribute:} \texttt{admissionDay : int}
        \item \textbf{Link:} \texttt{patientId[1..1] -> Patient}
        \item \textbf{Link:} \texttt{roomId[1..1] -> Room}
        \item \textbf{Link:} \texttt{operationTheatreId[1..1] -> OperatingTheatre}
    \end{itemize}
    \textbf{Nuance:} the admission stores only the start day, so full-stay effects are handled through day-based support objects such as \texttt{RoomAvailability}.

    \item \clabel{Room.}\\
    \textbf{Purpose:} represents a recovery room used for post-operative stays.
    \begin{itemize}
        \item \textbf{Attribute:} \texttt{id : int}
        \item \textbf{Attribute:} \texttt{maxCapacity : int}
    \end{itemize}

    \item \clabel{OperatingTheatre.}\\
    \textbf{Purpose:} represents a theatre resource for surgery.
    \begin{itemize}
        \item \textbf{Attribute:} \texttt{id : int}
    \end{itemize}

    \item \clabel{Surgeon.}\\
    \textbf{Purpose:} represents a surgeon resource.
    \begin{itemize}
        \item \textbf{Attribute:} \texttt{id : int}
    \end{itemize}
    \textbf{Nuance:} surgeons are assigned to patients before optimisation rather than chosen by the transformations.
\end{itemize}

\section*{Day-Based Resource State}
\begin{itemize}
    \item \clabel{SurgeonAvailability.}\\
    \textbf{Purpose:} stores how much operating time a surgeon has on a given day.
    \begin{itemize}
        \item \textbf{Attribute:} \texttt{day : int}
        \item \textbf{Attribute:} \texttt{maxOperatingTime : int}
        \item \textbf{Link:} \texttt{surgeonId[1..1] -> Surgeon}
    \end{itemize}

    \item \clabel{OperatingTheatreAvailability.}\\
    \textbf{Purpose:} stores available theatre capacity for one theatre on one day.
    \begin{itemize}
        \item \textbf{Attribute:} \texttt{day : int}
        \item \textbf{Attribute:} \texttt{maxCapacity : int}
        \item \textbf{Link:} \texttt{operatingTheatreId[1..1] -> OperatingTheatre}
    \end{itemize}

    \item \clabel{RoomAvailability.}\\
    \textbf{Purpose:} stores room occupancy state for one room on one day.
    \begin{itemize}
        \item \textbf{Attribute:} \texttt{day : int}
        \item \textbf{Attribute:} \texttt{occupiedBeds : int}
        \item \textbf{Attribute:} \texttt{ageGroup : AgeGroup}
        \item \textbf{Attribute:} \texttt{roomNumber : int}
        \item \textbf{Link:} \texttt{roomId[1..1] -> Room}
    \end{itemize}
    \textbf{Nuance:} this class is central to enforcing room capacity and no-age-group-mixing constraints across a patient’s full stay. The room number is duplicated here as a convenience field for transformation logic.
\end{itemize}

\section*{Nurse Assignment Layer}
\begin{itemize}
    \item \clabel{HospitalisationShift.}\\
    \textbf{Purpose:} represents nursing demand for a particular room on a particular day and shift.
    \begin{itemize}
        \item \textbf{Attribute:} \texttt{day : int}
        \item \textbf{Attribute:} \texttt{shift : int}
        \item \textbf{Attribute:} \texttt{roomNumber : int}
        \item \textbf{Attribute:} \texttt{workload : int}
        \item \textbf{Attribute:} \texttt{skillLevelRequired : int}
        \item \textbf{Link:} \texttt{room[1..1] -> Room}
        \item \textbf{Link:} \texttt{patient[0..*] -> Patient}
    \end{itemize}
    \textbf{Nuance:} this is the bridge between patient scheduling and nurse assignment because it expresses the demand that nurses must cover.

    \item \clabel{Nurse.}\\
    \textbf{Purpose:} represents a nurse resource.
    \begin{itemize}
        \item \textbf{Attribute:} \texttt{id : int}
        \item \textbf{Attribute:} \texttt{skillLevel : int}
    \end{itemize}

    \item \clabel{NurseWorkingShift.}\\
    \textbf{Purpose:} stores when a nurse is available to work and how much workload they can cover.
    \begin{itemize}
        \item \textbf{Attribute:} \texttt{day : int}
        \item \textbf{Attribute:} \texttt{shift : int}
        \item \textbf{Attribute:} \texttt{maxLoad : int}
        \item \textbf{Link:} \texttt{nurse[1..1] -> Nurse}
    \end{itemize}

    \item \clabel{RoomShiftAssignment.}\\
    \textbf{Purpose:} records the decision to assign a nurse to a particular hospitalisation shift.
    \begin{itemize}
        \item \textbf{Link:} \texttt{nurse[1..1] -> Nurse}
        \item \textbf{Link:} \texttt{hospitalisationShift[1..1] -> HospitalisationShift}
    \end{itemize}
    \textbf{Nuance:} this is the final decision object for the nurse-assignment optimisation stage.
\end{itemize}

\section*{Demand and Tracking Support}
\begin{itemize}
    \item \clabel{PatientDayDemand.}\\
    \textbf{Purpose:} describes how much nursing demand a patient contributes on a relative day of their stay.
    \begin{itemize}
        \item \textbf{Attribute:} \texttt{relativeDay : int}
        \item \textbf{Attribute:} \texttt{workloadProduced : int}
        \item \textbf{Attribute:} \texttt{skillLevelRequired : int}
        \item \textbf{Link:} \texttt{patient[1..1] -> Patient}
    \end{itemize}
    \textbf{Nuance:} this allows demand to depend on where the patient is in their stay rather than only on the absolute calendar day.

    \item \clabel{DeletedAdmissionsTracker.}\\
    \textbf{Purpose:} counts how many admission removals have been performed during search.
    \begin{itemize}
        \item \textbf{Attribute:} \texttt{count : int}
    \end{itemize}
    \textbf{Nuance:} this exists mainly to support the objective function that penalises removed patients.

    \item \clabel{Occupant.}\\
    \textbf{Purpose:} represents a patient-room occupancy relation with a stay length.
    \begin{itemize}
        \item \textbf{Attribute:} \texttt{stayLength : int}
        \item \textbf{Link:} \texttt{patientId[1..1] -> Patient}
        \item \textbf{Link:} \texttt{assignedRoomId[1..1] -> Room}
    \end{itemize}
    \textbf{Nuance:} in the current model, most practical occupancy reasoning is handled through \texttt{Admission} and \texttt{RoomAvailability}, so this class is more auxiliary than central.
\end{itemize}

\section*{Modelling Notes}
\begin{itemize}
    \item Patient scheduling is start-day based, but room state is maintained day by day through \texttt{RoomAvailability}.
    \item Nurse assignment is not done directly to patients; it is done to room-shifts through \texttt{HospitalisationShift} and \texttt{RoomShiftAssignment}.
    \item Several helper attributes such as \texttt{isScheduled}, \texttt{roomNumber}, and \texttt{AgeGroup.EMPTY} were introduced to make MDEO transformation logic easier to express.
    \item The metamodel already separates long-term entities such as patients and nurses from dynamic state such as availabilities, assignments, and deletion counters.
\end{itemize}

\end{document}
